% ============================================================
% Title page
% ============================================================
\begin{titlepage}
\centering
\vspace*{2cm}

{\LARGE \textbf{Evaluating Large Language Models for Personalised Credit Card Recommendation from User Spending Profiles}}

\vspace{2cm}

{\Large Hindvrat Singh Rao}

\vspace{1cm}

{\large A dissertation submitted in partial fulfilment of the requirements\\ for the degree of MSc Advanced Computer Science with Artificial Intelligence}

\vspace{2cm}

{\large Department of Computer and Information Sciences\\ University of Strathclyde}

\vspace{2cm}

{\large \today}

\vfill
\end{titlepage}

% ============================================================
% Declaration
% ============================================================
\chapter*{Declaration}
\addcontentsline{toc}{chapter}{Declaration}

This dissertation is submitted in part fulfilment of the requirements for the degree of MSc Advanced Computer Science with Artificial Intelligence of the University of Strathclyde.

I declare that this dissertation embodies the results of my own work and that it has been composed by myself. Following normal academic conventions, I have made due acknowledgement to the work of others.

This research did not involve human participants, real personal data, or any other activity requiring ethics approval, and therefore no application to the Departmental Ethics Committee was required (see Section~\ref{sec:ethics} for the full ethical-considerations statement).

I give permission to the University of Strathclyde, Department of Computer and Information Sciences, to provide copies of the dissertation, at cost, to those who may in the future request a copy of the dissertation for private study or research. \hfill (please tick) Yes \rule{0.8cm}{0.4pt} \; No \rule{0.8cm}{0.4pt}

I give permission to the University of Strathclyde, Department of Computer and Information Sciences, to place a copy of the dissertation in a publicly available archive. \hfill (please tick) Yes \rule{0.8cm}{0.4pt} \; No \rule{0.8cm}{0.4pt}

I declare that the word count for this dissertation (excluding title page, declaration, abstract, acknowledgements, table of contents, list of illustrations, references and appendices) is \textbf{9,479} words at the time of this draft, to be confirmed at final submission.

I confirm that I wish this to be assessed as a \textbf{Type 5} Dissertation.

\vspace{2cm}
\noindent Signed: \rule{5cm}{0.4pt} \hfill Date: \rule{4cm}{0.4pt}

% ============================================================
% Abstract
% ============================================================
\chapter*{Abstract}
\addcontentsline{toc}{chapter}{Abstract}

This dissertation evaluates whether a free-tier, open-weight Large Language Model (Meta's Llama 3.1, 8B parameters, accessed via the Groq API) can generate personalised UK credit-card recommendations competitive with a conventional content-based recommender, and how prompting strategy affects recommendation quality, factual reliability, explanation quality, and consistency. Using a fixed catalogue of 20 real UK credit cards and 300 synthetic user profiles --- spanning eight consumer archetypes with spending distributions calibrated against the ONS \textit{Living Costs and Food Survey 2022--23} --- three prompting strategies (zero-shot, structured, few-shot) were evaluated across 900 LLM inference calls, scored against a content-based cosine-similarity baseline using a four-criteria evaluation framework: recommendation-baseline overlap, deterministic factual correctness, LLM-as-judge explanation quality, and repeated-run consistency.

Prompting strategy was found to have a large, statistically significant effect on recommendation quality (Kruskal-Wallis $H=60.70$, $p=6.6\times10^{-14}$), but in a direction contrary to common expectations from the general prompt-engineering literature: zero-shot prompting significantly outperformed both structured and few-shot prompting on overlap with the baseline (mean overlap 0.687 vs.\ 0.517 and 0.539 respectively) and top-1 exact-match accuracy (0.307 vs.\ 0.077 and 0.023). However, structured prompting produced the most factually reliable responses (73.7\% of responses entirely free of checkable factual errors, against 55.0\% for zero-shot and 24.0\% for few-shot), demonstrating that recommendation accuracy and factual trustworthiness are distinct properties that a single-metric evaluation would fail to distinguish. Strikingly, an independent LLM judge rated few-shot prompting's explanations as the clearest and best-grounded of the three strategies (mean 2.907/3, against 2.730 for zero-shot and 2.650 for structured, $p < 0.0001$) despite few-shot being the least accurate and least factually reliable strategy on every other measure --- demonstrating that perceived explanation quality is not a reliable proxy for correctness. Zero-shot was also numerically the most stable strategy under repeated querying (mean pairwise Jaccard similarity 0.826, against 0.808 for few-shot and 0.776 for structured), though this difference did not reach statistical significance on the sample tested. No single prompting strategy dominates across all evaluated criteria.

These findings suggest that prompting-strategy choice for LLM-based recommendation should be validated empirically against the specific task and model rather than assumed from general NLP benchmarks, and that recommendation-quality and factual-reliability metrics should be reported and optimised for separately in applications where the reliability of the model's stated reasoning matters.

% ============================================================
% Acknowledgements
% ============================================================
\chapter*{Acknowledgements}
\addcontentsline{toc}{chapter}{Acknowledgements}

I would like to thank my dissertation supervisor, Dimitri Russinov, for their guidance and feedback throughout this project, particularly in shaping the baseline methodology and the scale of the experimental dataset.
